\documentclass[12pt,a4paper]{article}
\usepackage{graphicx}
\usepackage{ textcomp }
\setlength{\voffset}{-0.75in}
\setlength{\headsep}{5pt}

\begin{document}
	
	\begin{titlepage}
		\centering
		\includegraphics[width=0.4\textwidth]{./images/logo_ciencias}\par\
		{\LARGE\bf Facultad de Ciencias \par}
		{\LARGE\bf UNAM \par}
		\vspace{1cm}
		{\huge\bfseries Metodolog\'ia de desarrollo de software: Crystal Clear\par}
		\vspace{2cm}
		{\Large\itshape Oscar Eduardo Villa Chio\par}
		{\Large\itshape Jorge Luis Garc\'ia Flores\par}
		{\Large\itshape Enrique Antonio Bernal Cedillo\par}
		{\Large\itshape V\'ictor Zamora Guti\'errez\par}
		\vspace{0.5cm}
		{\large Riesgo Tecnol\'ogico\par 6\textdegree \ Semestre\par}
		% Bottom of the page
		{\large 22 de abril del 2016}
	\end{titlepage}
	
	\section{Introducci\'on}
	
	\subsection{Objetivo general}
	El objetivo general de este trabajo es introducir la metodolog\'ia Crystal Clear y explicar sus usos a trav\'es de varios ejemplos.
	
	\subsection{Objetivos espec\'ificos}
	\begin{itemize}
		\item Mostrarle a los alumnos de la clase de riesgo tecnol\'ogico c\'omo trabajar con la metodolog\'ia Crystal Clear.
		
		\item Otorgar una herramienta para terminar trabajos de peque\~no alcance en tiempos razonables.
	\end{itemize}
	
	\subsection{Alcance del trabajo}
	Clase de Riesgo Tecnol\'ogico.
	\pagebreak
	
	\section{Contenido del trabajo}
	Crystal Clear A Human-Powered Methodology For Small Teams, including The Seven Properties of Effective Software Projects \linebreak
	
	La metodolog\'ia \textit{Crystal Clear} fue desarrollada por Alistair Cockburn.\\
	En su libro, Alistair explica que dicha metodolog\'ia fue desarrollada para la elaboraci\'on de proyectos peque\~nos (de un aproximado de 8 integrantes como máximo, dentro de un area en com\'un). \\
	
	Para obtener esta metodolog\'ia, Alistair realiz\'o investigaci\'on sobre grupos peque\~nos de desarrollo que cumplian con \'exito lo requerido.\\
	Fue de esta manera que se encontr\'o que varios equipos realizaban las mismas acciones:	
	\begin{itemize}
	\item Sentar a la gente cerca. Que se comuniquen seguido y en un ambiente de amabilidad y de buena gana.
	
	\item Alejar los temas burocr\'aticos lo m\'as posible de los desarrolladores.
	
 	\item Dejar que los desarrolladores dise\~nen el sistema y tomen las decisiones que consideren apropiadas.
	
	\item Tener a un usuario involucrado en el proyecto.\\
	(Este ayuda a encontrar errores, realizar pruebas y perfeccionar el dise\~no del sistema)
	
	\item Tener un conjunto de pruebas de regresi\'on automatizadas para asegurarse de que el software funciona correctamente en cada paso del desarrollo.
	
	\item Producir funcionalidad entregable desde temprano y seguido (?).
	\\
	\end{itemize}
	
	La metodolog\'ia Crystal Clear se puede describir como sigue:
	
	El dise\~nador principal junto con otras 2 a 7 personas trabajan en un espacio peque\~no como una habitaci\'on (o varias habitaciones adyacentes) con muestras visuales de informaci\'on (como pizarrones) sin ninguna distracci\'on y con acceso directo a usuarios clave.\\
	
	Deben entregar c\'odigo usable y probado aproximadamente cada 3 meses.
	La metodolog\'ia Crystal Clear no pretende ser la mejor, sino tan s\'olo ser "suficiente".
	
	\section{Conclusiones}
	
	\section{Bibliografía}
	
	\section{Anexos}
	(Solo dejar en caso de requerirlos)

\end{document}